% ============================================================================ 6 · Información complementaria + conclusión
\section{Información complementaria}
\label{sec:complementaria}

\textbf{Documentación utilizada.} El conjunto de datos y su diccionario se recibieron en dos entregas, y el diccionario se
contrastó con los archivos columna por columna (\S\ref{sec:datos}). Cada hallazgo, decisión y experimento quedó documentado
con su fecha, su motivo y el procedimiento que lo reproduce, y cada experimento está definido por un archivo de configuración
versionado.

\textbf{Información que solo Ford puede aportar.} El límite actual no es el modelo sino la información disponible:
\begin{enumerate}
  \item \textbf{La tasa real de eventos} en la flota, que permite separar la física del criterio con que se extrajeron las
    listas. Es la principal incertidumbre del trabajo.
  \item \textbf{El origen de los 278 vehículos con evento y fecha por defecto} que salieron entre entregas, y el criterio
    de armado de la lista de vehículos con evento.
  \item \textbf{Qué registra la fecha del evento}: hoy se corre 21 días porque cae unas dos semanas después del taller.
  \item \textbf{Más eventos registrados} (hoy son 177), \textbf{la zona de uso} del vehículo (no la de venta) y \textbf{los
    costos reales} de una falla y de una inspección.
  \item \textbf{Un piloto en sombra de tres meses en un país}, para medir prevalencia, falsas alarmas y detección reales.
\end{enumerate}

\textbf{Trabajo futuro.} Reentrenar con los 557 vehículos y empaquetar el modelo. Validar toda mejora con datos nuevos (una
extracción posterior al 14-09-2026) y no con este conjunto de prueba, que ya fue utilizado. Medir en el piloto lo que estos
datos no permiten: si los avisos cambian hábitos, cuántas fallas se evitan y cuánto combustible se ahorra. Extender el
circuito a otras fallas del sistema de postratamiento.

% ============================================================================ Conclusión
\section{Conclusión}
\label{sec:conclusion}

Se mostró que la telemetría que Ford ya recibe contiene información que anticipa la saturación del filtro de partículas, y
que esa información es consistente con la física de la regeneración: los vehículos que fallan pasan más tiempo detenidos con
el motor en marcha y con el motor por debajo de su temperatura de régimen, y la diferencia crece hacia el evento. Se mostró
también que los datos contenían sesgos de muestreo y de formato capaces de producir métricas favorables sin ninguna
capacidad predictiva, y que el PR-AUC por fila tiene un piso de $0{,}177$ alcanzable sin información temporal. Por eso el
modelo se seleccionó con la detección por vehículo a presupuesto de falsas alarmas, comparada con un nulo que conserva el
historial.

Con ese criterio, el modelo finalista, una GRU con atención de $\sim 3.900$ parámetros en un ensamble de tres semillas,
detecta en vehículos no vistos $\sim 30\,\%$ de los eventos con $5\,\%$ de falsas alarmas y $\sim 50$--$60\,\%$ con $20\,\%$,
con una mediana de $\sim 4.600$\,km (unos tres meses) de anticipación y un umbral que respeta su presupuesto fuera de
muestra. Duplica la detección de los modelos tabulares y de supervivencia en el mismo conjunto. El circuito propuesto
convierte cada alerta en un aviso verificable, con el punto de operación como decisión de Ford, a un costo estimado de
USD~$0{,}07$--$0{,}10$ por vehículo y por año. La validación de su impacto sobre fallas evitadas requiere el piloto en sombra
propuesto.
